% ==========================================================================
%  Section 3 -- The lattice free field and its physical-volume law.  Owner: W03.
%  Frozen labels (CONTRACT section 4): sec:lattice, def:lattice-laplacian,
%  prop:operator, def:gj-measure, def:finite-source-lift, rem:lattice-input,
%  lem:haar-lebesgue; added here: lem:finite-parseval and the eq: labels.
%
%  Proved in this file: the eigenfunction identity and the positivity bound
%  (prop:operator), finite Parseval (lem:finite-parseval), the Haar normalisation
%  (lem:haar-lebesgue).  Defined: the Gaussian operator and its symbol, the
%  physical lattice Gaussian law, the finite-source Fourier lift and its law.
%  The lattice covariance is an input (rem:lattice-input).  Nothing here states
%  a limit; the two limits are named in the closing paragraph.
% ==========================================================================
\section{The lattice free field and its physical-volume law}\label{sec:lattice}

This section fixes the two objects that Part~I consumes. \Cref{def:lattice-laplacian} introduces the
discrete Laplacian and the Gaussian operator; \cref{prop:operator} proves the eigenfunction identity and
the positivity bound that every later estimate uses. \Cref{lem:finite-parseval} is the one estimate of the
paper that spends boundedness, \cref{def:gj-measure} is the physical lattice Gaussian law, and
\cref{def:finite-source-lift} carries that law to the weighted Fourier side. \Cref{lem:haar-lebesgue}
reconciles the two Fourier normalisations of \cref{rem:v-normalisation}. Two points made in passing matter
later: the lattice covariance is an input and not derived here (\cref{rem:lattice-input}), and the spacing
enters more than once, in different powers (\cref{def:gj-measure}).

\subsection{The discrete Laplacian and the Gaussian operator}

\begin{definition}[The discrete Laplacian and the Gaussian operator]\label{def:lattice-laplacian}
Let $N\ge1$ and $a=L/N$ as in \cref{hyp:volume}. A \emph{configuration} is a real function $\omega$ on the
lattice $I_N=(\Z/N\Z)^2$ of \cref{def:modes}; the lattice configuration space is $\R^{I_N}$, and we also
admit complex $u:I_N\to\C$. With $e_1,e_2$ the unit residues, the \emph{discrete Laplacian} is
\begin{equation}\label{eq:laplacian}
  (\Delta_au)(x):=a^{-2}\sum_{i=1}^{2}\bigl(u(x+e_i)+u(x-e_i)-2u(x)\bigr),
\end{equation}
the two oriented neighbours being subtracted and added exactly once, so that \eqref{eq:laplacian} keeps
its form even when $N\le2$. The \emph{Gaussian operator} is
\begin{equation}\label{eq:gaussian-operator}
  Q:=m^2-\Delta_a ,
\end{equation}
and we write $A_N:=Q$ when $Q$ acts on configurations against the \emph{counting} inner product, the one
attached to the vector space $\R^{I_N}$:
\begin{equation}\label{eq:counting-inner-product}
  \ip{u}{v}_{\mathrm{counting}}:=\sum_{x\in I_N}u(x)v(x),
\end{equation}
so that $A_N^{-1}$ is the inverse of $Q$. For $p\in I_N$ the \emph{character} at $p$ is the function
\[
  \chi_p(x):=\exp\bigl(2\pi i\,p\cdot x/N\bigr),\qquad x\in I_N ,
\]
the product $p\cdot x$ being read modulo $N$.
\end{definition}

\begin{proposition}[Operator facts]\label{prop:operator}
$\Delta_a$ is a symmetric real matrix, hence self-adjoint, and $Q$ is positive definite with spectrum
contained in $[m^2,\infty)$. Consequently $Q$ is invertible and has a positive square root $Q^{1/2}$, and
$\C^{I_N}$ has an orthonormal basis of eigenvectors of $Q$. The eigenfunctions are the characters, with
eigenvalues the symbol $\lambda_N$ of \eqref{eq:lattice-symbol}:
\begin{equation}\label{eq:eigen-identity}
  \Delta_a\chi_p=-4a^{-2}\Bigl(\sum_{i=1}^{2}\sin^2\frac{\pi p_i}{N}\Bigr)\chi_p ,
  \qquad
  Q\chi_p=\lambda_N(p)\,\chi_p ,
\end{equation}
and $\lambda_N(p)\ge m^2>0$ for every $p\in I_N$.
\end{proposition}

\begin{proof}
\emph{The eigenfunction identity.} Translating the argument of a character multiplies it by a scalar: for
$i=1,2$ and every $x$,
\[
  \chi_p(x+e_i)=\zeta_i\chi_p(x),\qquad \chi_p(x-e_i)=\zeta_i^{-1}\chi_p(x),
  \qquad \zeta_i:=e^{2\pi ip_i/N},
\]
because $p_i(x_i+1)\equiv p_ix_i+p_i$ modulo $N$. Inserting this into \eqref{eq:laplacian} gives
$(\Delta_a\chi_p)(x)=a^{-2}\sum_i(\zeta_i+\zeta_i^{-1}-2)\chi_p(x)$, and Euler's formula gives
$\zeta_i+\zeta_i^{-1}-2=2\cos(2\pi p_i/N)-2=-4\sin^2(\pi p_i/N)$. Hence the eigenfunction identity of
\eqref{eq:eigen-identity}; the relation $Q\chi_p=\lambda_N(p)\chi_p$ then follows from the definition
\eqref{eq:lattice-symbol} of the symbol together with $Q=m^2-\Delta_a$.

\emph{Positivity.} Each summand $4a^{-2}\sin^2(\pi p_i/N)$ of \eqref{eq:lattice-symbol} is a nonnegative
real number, so $\lambda_N(p)\ge m^2$; and $m>0$ by \cref{hyp:volume}.

\emph{An explicit orthonormal eigenbasis.} Summing the one-coordinate geometric series and using
$\zeta^N=1$ for $\zeta=e^{2\pi i/N}$,
\[
  \sum_{x\in I_N}\chi_{-p}(x)\chi_q(x)
  =\prod_{i=1}^{2}\sum_{x_i\in\Z/N\Z}e^{2\pi i(q_i-p_i)x_i/N}=N^2\mathbf 1_{p=q} ,
\]
so the $N^2$ characters are pairwise orthogonal and nonzero, hence a basis of $\C^{I_N}$; scaled by
$N^{-1}$ they are orthonormal, and \eqref{eq:eigen-identity} makes each of them an eigenvector.

\emph{Self-adjointness and the spectrum.} The coefficient of $u(y)$ in $(\Delta_au)(x)$ is $-2a^{-2}$ when
$y=x$, $a^{-2}$ when $y-x=\pm e_i$ for one $i$, and $0$ otherwise; this is symmetric in $x$ and $y$, so
$\Delta_a$ is a symmetric real matrix, hence self-adjoint, and $Q$ is too. On the basis just exhibited $Q$
acts diagonally with the real entries $\lambda_N(p)$, so its spectrum is exactly
$\{\lambda_N(p):p\in I_N\}\subseteq[m^2,\infty)$: the operator is positive definite, hence invertible. A
positive definite symmetric matrix has a positive definite symmetric square root, and that is the
finite-dimensional spectral theorem, the only fact named in this proof; the eigenbasis itself is
exhibited above rather than postulated. Write $Q^{1/2}$ for that root.
\end{proof}

\begin{lemma}[Finite Parseval]\label{lem:finite-parseval}
Let $f,g:\R^2\to\R$ be bounded, and for $p\in I_N$ let
\begin{equation}\label{eq:sampled-coefficient}
  F_N(p):=a^2\sum_{x\in I_N}f(ax)\,\chi_{-p}(x),\qquad
  G_N(p):=a^2\sum_{x\in I_N}g(ax)\,\chi_{-p}(x)
\end{equation}
be the sampled coefficients at the physical points $ax$ of \cref{hyp:volume}, the factor $a^2$ being the
area of one cell. Then
\begin{equation}\label{eq:parseval}
  \sum_{p\in I_N}\overline{F_N(p)}\,G_N(p)=Va^2\sum_{x\in I_N}f(ax)g(ax),
  \qquad
  \sum_{p\in I_N}\abs{F_N(p)}^2=Va^2\sum_{x\in I_N}\abs{f(ax)}^2\le V^2\norm{f}_\infty^2 .
\end{equation}
\end{lemma}

\begin{proof}
Summing the one-coordinate geometric series over modes instead of over sites gives the dual orthogonality
relation $\sum_{p\in I_N}\chi_p(x)\chi_{-p}(y)=N^2$ when $x=y$ and $0$ otherwise. Since $f$ is real,
$\overline{\chi_{-p}(x)}=\chi_p(x)$; inserting \eqref{eq:sampled-coefficient} and interchanging the finite
sums therefore collapses the cross terms onto the diagonal:
\[
  \sum_{p\in I_N}\overline{F_N(p)}\,G_N(p)
  =a^4\!\!\sum_{x,y\in I_N}\!\!f(ax)g(ay)\sum_{p\in I_N}\chi_p(x)\chi_{-p}(y)
  =a^4N^2\sum_{x\in I_N}f(ax)g(ax).
\]
The bilinear identity of \eqref{eq:parseval} follows from $a^2N^2=a^2\cdot L^2/a^2=V$; the diagonal one is
the case $g=f$, and $a^2\sum_x\abs{f(ax)}^2\le a^2N^2\norm{f}_\infty^2=V\norm{f}_\infty^2$.
\end{proof}

This is the only place in Part~I where a test is spent through boundedness rather than through smoothness:
the bound is $\norm{f}_\infty$ alone, and that is what lets \cref{sec:covariance} work with bounded
continuous periodic tests and ask for no derivative.

\subsection{The physical lattice Gaussian law}

\begin{definition}[The physical lattice Gaussian law]\label{def:gj-measure}
With $A_N=Q$ of \eqref{eq:gaussian-operator}, put
\begin{equation}\label{eq:gj-covariance}
  C_N:=a^{-1}A_N^{-1/2},\qquad\text{so that}\qquad C_N^2=a^{-2}A_N^{-1};
\end{equation}
the squared object $C_N^2$ is the covariance operator. The \emph{physical lattice Gaussian law}
$\mu^{\mathrm{GJ}}_N$ is the centred Gaussian Borel probability measure on $\R^{I_N}$ that is the law of
$C_N\xi$ for a standard Gaussian vector $\xi$ on $\R^{I_N}$, so that $C_N^2$ is its covariance operator. Its
test averages are
$\omega\mapsto\omega(f):=\sum_{x\in I_N}a^2f(ax)\,\omega(x)$, so that
$\omega(f)=\ip{b}{\omega}_{\mathrm{counting}}$ with the sampled vector $b_x:=a^2f(ax)$, $x\in I_N$. For
$b,c\in\R^{I_N}$ the covariance pairing is therefore
\begin{equation}\label{eq:gj-pairing}
  \ip{C_Nb}{C_Nc}=a^{-2}\ip{b}{A_N^{-1}c}_{\mathrm{counting}} .
\end{equation}
It is the factor $C_N$, not the operator $C_N^2$, that the rest of Part~I calls the lattice covariance.
\end{definition}

Only two properties of $\mu^{\mathrm{GJ}}_N$ are used later, and they are its operative content here. The
mean is exactly zero: $\int\omega(f)\,\mu^{\mathrm{GJ}}_N(\dd\omega)=0$ for every bounded $f$, because
$\omega(f)$ is linear in $\omega$ and the law is centred. The covariance is the inverse-mass identity
\eqref{eq:gj-pairing} in sampled form: with $b_x=a^2f(ax)$ and $c_x=a^2g(ax)$,
\begin{equation}\label{eq:sampled-covariance}
  \operatorname{Cov}\bigl(\omega(f),\omega(g)\bigr)
  =\ip{C_Nb}{C_Nc}
  =a^{-2}\ip{b}{A_N^{-1}c}_{\mathrm{counting}}
  =a^2\ip{f}{A_N^{-1}g}_{\mathrm{counting}} ,
\end{equation}
the last equality because the counting sum picks up one factor $a^2$ per sampled vector. No other
property of $\mu^{\mathrm{GJ}}_N$ enters any estimate below.

Where the powers of $a$ sit is worth stating once, explicitly, because the factors pull in opposite
directions and a reader who cancels them gets a constant wrong. There are three occurrences and they have
three different jobs. \emph{(i)} The cell area $a^2$ sits on the \emph{test evaluation}: it is the
quadrature weight of the sum $\omega(f)=\sum_xa^2f(ax)\omega(x)$, hence part of the sampled vector
$b_x=a^2f(ax)$ in \eqref{eq:sampled-covariance}. \emph{(ii)} The covariance factor carries
$(a^2)^{-1/2}=a^{-1}$, and the pairing \eqref{eq:gj-pairing} squares it, so $a^{-2}$ sits on the bilinear
form. \emph{(iii)} The
same $a^2$ reappears inside the per-site vector of the lift \eqref{eq:lift}, where it is again the cell
area and is again a quadrature weight, this time a Fourier one. The cancellation between \emph{(i)} and
\emph{(ii)} is exact and leaves the single factor $a^2$ of \eqref{eq:sampled-covariance}, since
$a^{-2}(a^2)^2=a^2$. The $a^2$ of \emph{(iii)} is not the inverse of the $a^{-2}$ of \emph{(ii)}: the two
belong to different functionals, and neither may be cancelled against the other.

\begin{remark}[The lattice covariance is an input]\label{rem:lattice-input}
$C_N$ is data: \eqref{eq:gj-covariance} fixes it, and \cref{prop:operator} is what makes that
specification consistent, since $Q$ is then invertible with a positive square root. It is not derived here
from any field-theoretic construction, and no property of $C_N$ beyond exact zero mean and
\eqref{eq:gj-pairing} is used anywhere below. The name GJ attached to this normalisation records which
power of the spacing is meant --- namely $a^{-1}$ per factor in \eqref{eq:gj-covariance} --- and it is a
name and not a citation: no step of the proof is taken from a reference.
\end{remark}

\subsection{The finite-source Fourier lift and its induced law}

\begin{definition}[The finite-source Fourier lift]\label{def:finite-source-lift}
For a configuration $\omega\in\R^{I_N}$ and a mode $p\in\Z^2$ put
\begin{equation}\label{eq:lift}
  (L_N\omega)(p):=\sqrt{w_3(p)}\;a^2\sum_{x\in I_N}\omega(x)\,e^{-2\pi i\,p\cdot(ax)/L},
\end{equation}
with $w_3$ the weight of \cref{def:modes}, and let
\begin{equation}\label{eq:finite-source-law}
  \mu_N:=(L_N)_*\,\mu^{\mathrm{GJ}}_N
\end{equation}
be the image of the law of \cref{def:gj-measure} under $L_N$.
\end{definition}

The phase in \eqref{eq:lift} is the physical one, $e^{-2\pi i\,p\cdot(ax)/L}$, and since $a/L=1/N$ it
equals $\chi_{-p}(x)$: the lift is the discrete Fourier transform of the sampled field, not of the residue
index, and that is what makes $a^2$ the cell area rather than a bookkeeping symbol. Because the phase has
period $N$ in each component of $p$, the family \eqref{eq:lift} is the $N^2$-periodic repetition over
$\Z^2$ of its values on $I_N$. Each entry is a finite real-linear combination of the coordinate
evaluations $\omega\mapsto\omega(x)$, so $L_N$ is Borel as a map into $\C^{\Z^2}$; and $L_N$ lands in the
carrier of \cref{def:carrier-hs}, not merely in some larger space, because with
$\norm{\omega}_1:=\sum_x\abs{\omega(x)}$,
\[
  \sum_{p\in\Z^2}\abs{(L_N\omega)(p)}^2
  \le a^4\norm{\omega}_1^2\sum_{p\in\Z^2}w_3(p)
  =a^4\norm{\omega}_1^2\Bigl(1+\sum_{p\neq0}\norm{p}_\infty^{-6}\Bigr)<\infty
\]
by \cref{ext:summability}, and $w_3\le1$ makes every weighted sum of \cref{def:carrier-hs} no larger than
this one. So \eqref{eq:finite-source-law} is the genuine pushforward of a Borel probability measure by a
Borel map into the carrier: $\mu_N(A)=\mu^{\mathrm{GJ}}_N(L_N^{-1}A)$ for Borel $A$, and $\mu_N$ is a
Borel probability measure there for every $N$.

What \eqref{eq:finite-source-law} \emph{is} is a carrier-valued representative of the lattice measure
itself, and \cref{lem:lift-representative} records this as a lemma: $L_N$ is total on the
configuration space, Borel, and lands in the carrier. What is \emph{not} claimed anywhere is a
statement about a total lift specified independently of $L_N$ --- in particular the series-defined
total lift of the machine development, where the measurability of that chosen function is an input
of the data structure rather than a construction (\cref{rem:lift-development}). The object of the
law limit is $L_N$ and the measure $\mu_N$ of \eqref{eq:finite-source-law}.

\subsection{Haar versus physical Lebesgue normalisation}

\begin{lemma}[Haar-to-physical-Lebesgue normalisation]\label{lem:haar-lebesgue}
Let $f$ be a bounded continuous $L$-periodic test, put $f_L(z):=f(Lz)$ for $z\in[0,1)^2$, and let $\dd z$
be Lebesgue measure on the unit cell, of total mass $1$. The orthonormal coefficient of $f_L$ at
$k\in\Z^2$ is
\begin{equation}\label{eq:haar-coefficient}
  (Bf_L)(k):=\int_{[0,1)^2}f(Lz)\,e^{-2\pi i\,k\cdot z}\,\dd z=V^{-1}F(k),
\end{equation}
with $F$ the integral-normalised coefficient \eqref{eq:fourier-coefficient} of \cref{def:torus}.
\end{lemma}

\begin{proof}
Substituting $x=Lz$, so that $\dd x=L^2\dd z$ and $k\cdot z=k\cdot x/L$, three steps give the claim:
\begin{align*}
  (Bf_L)(k)
  &=\int_{(0,1]^2}f(Lz)\,e^{-2\pi i\,k\cdot z}\,\dd z ,\\
  &=\int_{[0,L]^2}f(x)\,e^{-2\pi i\,k\cdot x/L}\,L^{-2}\,\dd x ,\\
  &=V^{-1}\int_{[0,L]^2}f(x)\,e^{-2\pi i\,k\cdot x/L}\,\dd x=V^{-1}F(k) .
\end{align*}
The half-open cell may be replaced by the closed one because the two differ by a Lebesgue-null set; the
phase identity $k\cdot(Lz)=L\,(k\cdot z)$ holds for negative $k$ as written, with no conjugation; and
$L^2=V$.
\end{proof}

This is the single easiest place in the paper to lose a factor of $L^d$, which in two dimensions is the
volume $V$: a change of variables pays exactly one Jacobian, and substituting the orthonormal coefficient
where the integral-normalised one is meant, inside a pairing that already carries its own $V^{-1}$
(\cref{rem:v-normalisation}), loses that factor if it is done once and injects a spurious one if it is done
twice.

\subsection{The two limits taken later}

Two limits are taken in what follows, and they are of different kinds. One is the covariance limit of
\cref{sec:covariance}: the lattice covariance pairings \eqref{eq:sampled-covariance} of bounded continuous
periodic tests are shown to converge as $N\to\infty$ at fixed $L$ and $m$, against an independently
constructed resolvent. The other is the law limit of \cref{sec:tightness}--\cref{sec:main}: the object is
the sequence $(\mu_N)$ of pushforwards \eqref{eq:finite-source-law}, each a Borel probability measure on
the carrier, and the question is what a limit of that sequence on that carrier is. Which topology carries
it, and what a limit there means, is fixed in \cref{sec:carrier}; the definitions above are what that
section and the ones after it take as given.
